% wave14_j4_chiral_bar_cobar_factorization.tex
% Wave-14 J4/20: explicit chiral bar--cobar $B^{\chir} \dashv \Omega^{\chir}$
% for 6D hCS factorisation algebras.
% Channel: Francis--Gaitsgory + Beilinson.
% Lifts Wave-13 G5 ($\CE^*_{\dbar,\chir} = E_3$-algebra) from \emph{object-
% level identification} to \emph{adjunction-level Quillen equivalence},
% closing the bar--cobar loop on the chiral side and bridging Vol~I Thm~A.

\providecommand{\Obs}{\mathrm{Obs}}
\providecommand{\hCS}{\mathrm{hCS}}
\providecommand{\CE}{\mathrm{CE}}
\providecommand{\Ch}{\mathrm{Ch}}
\providecommand{\Dolb}{\mathrm{Dolb}}
\providecommand{\chir}{\mathrm{chir}}
\providecommand{\dbar}{\bar\partial}
\providecommand{\kch}{\kappa_{\mathrm{ch}}}
\providecommand{\kcat}{\kappa_{\mathrm{cat}}}
\providecommand{\Ran}{\mathrm{Ran}}
\providecommand{\Sym}{\mathrm{Sym}}
\providecommand{\Hom}{\mathrm{Hom}}
\providecommand{\C}{\mathbb{C}}
\providecommand{\cO}{\mathcal{O}}
\providecommand{\cE}{\mathcal{E}}
\providecommand{\cF}{\mathcal{F}}
\providecommand{\cC}{\mathcal{C}}
\providecommand{\cL}{\mathcal{L}}
\providecommand{\g}{\mathfrak{g}}
\providecommand{\Lie}{\mathrm{Lie}}
\providecommand{\coLie}{\mathrm{coLie}}
\providecommand{\coCE}{\mathrm{coCE}}
\providecommand{\Bar}{\mathrm{B}}
\providecommand{\Cob}{\Omega}

\section*{Chiral bar--cobar $B^{\chir} \dashv \Omega^{\chir}$ for 6D hCS
  (Wave-14 J4/20)}

\paragraph{Theorem (main).}
On factorisation algebras over $\C^3$ with Dolbeault differential, the
chiral bar and cobar functors
\[
  B^{\chir}:\; \Lie^{\chir}(\Ch(\Dolb)) \;\rightleftarrows\;
                \coLie^{\chir}(\Ch(\Dolb)) \;:\! \Cob^{\chir}
\]
form a Quillen adjunction whose restriction to connected chiral Lie
algebras is a Quillen equivalence. Applied to $\cE_{\hCS} = \Omega^{0,
\bullet}(\C^3, \g)$,
\[
  B^{\chir}(\cE_{\hCS}) \;\simeq\; \CE^*_{\dbar,\chir}(\cE_{\hCS},\,\C),
  \qquad
  \Cob^{\chir}\!\bigl(\CE^*_{\dbar,\chir}(\cE_{\hCS},\,\C)\bigr) \;\simeq\;
  \cE_{\hCS}
\]
as chiral Lie algebras in $\Ch(\Dolb)$, closing the loop witnessed at
the enveloping level by Wave-13 G5's
$U^{\chir}(\cE_{\hCS}) \simeq \Omega_{E_3}(\CE^*_{\dbar,\chir})$.

\paragraph{Cycle 1: bar $B^{\chir}$, explicit.}
\emph{Attack.} Francis--Gaitsgory 2012 (Selecta Math.\ 18, \S3.1;
henceforth FG) define, for any chiral Lie algebra $\cF$ on a smooth
curve $X$, the chiral CE complex $\CE^*_{\chir}(\cF)$ as the
$\cD_X$-module
\[
  \CE^*_{\chir}(\cF) \;:=\; \Sym^\bullet_{\chir}\!\bigl(\cF[1]\bigr)^\vee,
  \qquad d_{\CE} = d_{\cF} + d_{\text{chir}},
\]
with $d_{\text{chir}}$ the chiral bracket (BD04 Ch.~3, Def.~3.4.11)
acting on the shuffle-symmetric tensors supported along diagonals
in $X^n$.
\emph{Heal for $\C^3$.} For the Dolbeault-enhanced setting of 6D hCS
we replace $\cD_X$ by Dolbeault-dressed $\cD^{\dbar}_{\C^3}$: the chiral
Sym on $\Conf_n(\C^3)$ is built from shuffle-symmetric tensors of
$(0,\bullet)$-forms along partial diagonals, with differential
$d_{\CE,\dbar} = \dbar + d_{\cF} + d_{\text{chir}}$. Explicitly for
$\cF = \cE_{\hCS}$,
\[
  B^{\chir}(\cE_{\hCS}) \;=\;
  \Bigl(\bigoplus_{n \geq 0} j_*j^* \Sym^n_{\chir}(\cE_{\hCS}[1])^\vee,
        \; \dbar + d_{\text{chir}}\Bigr)
  \;\simeq\; \CE^*_{\dbar,\chir}(\cE_{\hCS},\,\C),
\]
with $j: \Conf_n(\C^3) \hookrightarrow (\C^3)^n$ the open embedding
of configuration space. The $\Sym^n_{\chir}$ piece lives on the
Ran space $\Ran(\C^3)$; restricting to $n$-fold configurations gives
the diagonal-supported shuffle.

\paragraph{Cycle 2: cobar $\Cob^{\chir}$, explicit.}
\emph{Attack.} FG Thm~3.3.1: on $\coLie^{\chir}$, the cobar functor
$\Cob^{\chir}$ is chiral CE-chains of the free chiral Lie resolution,
explicitly
\[
  \Cob^{\chir}(\cC) \;:=\; \Bigl(\bigoplus_{n \geq 1} j_*j^*
    \Sym^n_{\chir}\!\bigl(\cC[-1]\bigr), \; \partial_{\cC} + \partial_{\text{chir}}\Bigr)
\]
with $\partial_{\text{chir}}$ the chiral cobracket acting as the dual
of the chiral bracket on $B^{\chir}$. The grading shift $\cC[-1]$
inverts the shift $\cF[1]$ of Cycle~1, and the differential dualises
diagonal support to meromorphic poles.
\emph{Heal.} For connected chiral Lie coalgebras
$\cC \in \coLie^{\chir}_{\text{conn}}(\Ch(\Dolb))$ --- those with
$\cC^{(0)} = 0$ so the bar filtration is complete --- the
Mittag--Leffler condition on the chiral filtration $F^p \cC =
\bigoplus_{n \geq p} \cC^{(n)}$ forces convergence of the cobar
spectral sequence (BD04 \S3.4.17). Without connectedness the cobar
diverges on the chiral completion; FG \S3.3 localises at connected
coalgebras exactly to avoid this.

\paragraph{Cycle 3: Quillen equivalence on connected parts.}
\emph{Attack.} FG Thm~6.2.1 (the core theorem of chiral Koszul
duality) states: the bar--cobar adjunction
\[
  B^{\chir}:\; \Lie^{\chir}_{\text{conn}}(\Ch(X))
  \;\rightleftarrows\;
  \coLie^{\chir}_{\text{conn}}(\Ch(X)) \;:\! \Cob^{\chir}
\]
is a Quillen equivalence for $X$ a smooth curve, where model
structures are the chiral-projective structures (FG \S2.4).
\emph{Heal for $\Ch(\Dolb)$.} On $\C^3$ we replace $X$ by the
Dolbeault-dressed complex manifold and the curve-chiral operad by
the holomorphic-factorisation operad of Gwilliam--Williams 2021
(Thm~2.5.5: holomorphic factorisation algebras on $\C^d
\leftrightarrow E_d^{\mathrm{hol}}$-algebras). The FG proof extends
by two inputs: (i) the Dolbeault resolution
$\Omega^{0,\bullet}(\C^3) \to \cO^{\mathrm{hol}}_{\C^3}$ is a
cofibrant replacement in $\Ch(\Dolb)$, so chiral-projective
replacement lifts from the algebraic to the Dolbeault setting;
(ii) the configuration-space Fadell--Neuwirth fibration
$\Conf_{n+1}(\C^3) \to \Conf_n(\C^3)$ has fibre $\C^3 \setminus
\{n \text{ points}\}$ which is $(2 \cdot 3 - 2) = 4$-connected,
so the higher chiral-Sym pieces contribute no new obstructions
beyond the algebraic case. The unit $\mathrm{id} \to \Cob^{\chir}
B^{\chir}$ and counit $B^{\chir} \Cob^{\chir} \to \mathrm{id}$ are
quasi-isomorphisms on connected objects by direct bar-filtration
spectral-sequence computation (BD04 Prop.~3.4.19 transcribed to
Dolbeault differentials).

\paragraph{Cycle 4: bar of hCS is chiral CE with coefficients in $\C$.}
\emph{Attack.} Apply $B^{\chir}$ to $\cE_{\hCS}$.
\emph{Construction.} The chiral Lie structure on $\cE_{\hCS} =
\Omega^{0,\bullet}(\C^3, \g)$ is the pointwise $\g$-bracket
extended chirally to Ran-space via diagonal support (Costello--
Gwilliam 2021 Vol~II, \S5.4). Its bar is, by Cycle~1,
\[
  B^{\chir}(\cE_{\hCS}) \;=\; \CE^*_{\dbar,\chir}(\cE_{\hCS},\,\C),
\]
the $\C$-coefficient chiral CE cochain complex (Wave-13 E4
Def.~\ref{def:dolbeault-chiral-ce}, specialised to the trivial
$\cE_{\hCS}$-module $\C$). \emph{Heal.} Wave-13 G5 established
the $E_3$-algebra equivalence
\[
  \Obs_{\hCS}(\C^3) \;\simeq\;
  \CE^*_{\dbar,\chir}(\cE_{\hCS},\,\cO_{\C^3})
\]
as observable algebras. Cycle~4 refines this: taking coefficients
in $\C$ (rather than $\cO_{\C^3}$) identifies the bar of the
\emph{Lie algebra} $\cE_{\hCS}$, which is the universal cochain
carrying the chiral-Lie information free of its module structure.
The relation is
\[
  \CE^*_{\dbar,\chir}(\cE_{\hCS},\,\cO_{\C^3}) \;=\;
   B^{\chir}(\cE_{\hCS}) \otimes_{\C} \cO_{\C^3}
\]
(base change along $\C \hookrightarrow \cO_{\C^3}$), consistent with
FG \S5 on CE-coefficient functoriality.

\paragraph{Cycle 5: reconstruction $\Cob^{\chir} B^{\chir}(\cE_{\hCS})
  \simeq \cE_{\hCS}$.}
\emph{Attack.} The counit
$\varepsilon: \Cob^{\chir} B^{\chir}(\cE_{\hCS}) \to \cE_{\hCS}$ is the
structure map of the Quillen equivalence on connected parts
(Cycle~3). \emph{Heal.} $\cE_{\hCS}$ is \emph{connected}: as a
chiral Lie algebra it is concentrated in non-negative Sym-degree
with $\cE_{\hCS}^{(0)} = 0$ because $\Omega^{0,\bullet}(\C^3, \g)$
starts at Sym-degree one. Projective-cofibrant replacement is
provided by the Koszul resolution
$P^\bullet \twoheadrightarrow \cE_{\hCS}$ where $P^\bullet = \Cob^{\chir}
\Sigma^{-1} \cE_{\hCS}^{\vee,\chir}$ is the chiral cobar of the
degree-shifted dual coalgebra (FG Prop.~4.1.7 adapted via
Gwilliam--Williams 2021). On $P^\bullet$ the counit is a
quasi-isomorphism by bar--cobar telescoping; transporting along
$P^\bullet \to \cE_{\hCS}$ gives
\[
  \Cob^{\chir}\!\bigl(\CE^*_{\dbar,\chir}(\cE_{\hCS},\,\C)\bigr) \;\simeq\;
  \cE_{\hCS} \qquad \text{in } \Lie^{\chir}(\Ch(\Dolb)).
\]

\paragraph{Bridge to Volume I Theorem A.}
Vol~I Thm~A (chiral bar--cobar) states the same adjunction on
curves; Wave-14 J4 is its $d = 3$ extension. The bridge map: at
$d = 1$ the enveloping algebra $U^{\chir}(\cF)$ is an $E_2$-chiral
algebra (vertex algebra after completion); at $d = 3$, by Wave-13 G5,
$U^{\chir}(\cE_{\hCS}) \simeq \Omega_{E_3}(B^{\chir}(\cE_{\hCS}))$ is
$E_3$-chiral on $\Ch(\Dolb)$. The two statements agree after chain-
level localisation: the FG curve-adjunction transports to the
Dolbeault-dressed $\C^3$-adjunction by holomorphic-operad dimension
shift ($n \to n + 2$ on $E$-degree, from $E_1$ to $E_3$). This is
CY-D stratification (CLAUDE.md): $\Phi$ outputs $E_1$ at $d \geq 3$
for $A = H^0\Obs$ while the bar--cobar on factorisation algebras
retains the full $E_3$-chiral content before passing to cohomology.

\paragraph{Consequence for $\kch$.} The closed loop gives, for compact
CY$_3$ $X$ with abelian $\g$,
\[
  \kch(A_X) \;=\; \chi\!\bigl(B^{\chir}(\cE_X)\bigr) \;=\;
   \sum_{q} (-1)^q h^{0,q}(X),
\]
the Hodge-supertrace on Dolbeault cohomology, matching the CLAUDE.md
essential constant. For the quintic ($h^{0,0} = 1, h^{0,1} = 0,
h^{0,2} = 0, h^{0,3} = 1$): $\kch = 2$. The computation is on
$B^{\chir}$ directly, no $\Omega^{\chir}$ required at the Euler-char
level.

\paragraph{Scope.}
\emph{Chain-level}: explicit shuffle sums, Dolbeault dressing of
chiral Sym, configuration-space connectivity bound, Koszul resolution
$P^\bullet$. \emph{$(\infty,1)$-categorical}: Quillen equivalence in
$\Lie^{\chir}_{\text{conn}}(\Ch(\Dolb))$ via FG Thm~6.2.1 lifted by
Gwilliam--Williams 2021 Thm~2.5.5. Both lanes load-bearing (CLAUDE.md
equal-status rule). The adjunction is \emph{structural}: Vol~I
curve-level, Wave-13 G5 object-level $E_3$, Wave-14 J4 functor-level
$B^{\chir} \dashv \Cob^{\chir}$ on $\C^3$.

\paragraph{References.}
\begin{itemize}
  \item J.\ Francis and D.\ Gaitsgory, ``Chiral Koszul duality,''
    \emph{Selecta Math.} 18 (2012), 27--87; Thm~3.3.1 (cobar),
    Thm~6.2.1 (Quillen equivalence), \S2.4 (model structure),
    Prop.~4.1.7 (Koszul resolution). arXiv:1103.5803.
  \item A.\ Beilinson and V.\ Drinfeld, \emph{Chiral Algebras},
    AMS Colloquium Publ.\ 51, 2004, Ch.~3 (chiral Lie algebras and
    chiral CE), Def.~3.4.11, Prop.~3.4.19.
  \item J.\ Lurie, \emph{Higher Algebra}, \S5.2 ($E_n$-algebras);
    \S5.5.3 (locally constant factorisation algebras).
  \item O.\ Gwilliam and B.\ R.\ Williams, ``Higher Kac--Moody
    algebras and symmetries of holomorphic field theories,''
    arXiv:2009.05037, Thm~2.5.5 (holomorphic factorisation on $\C^d
    \leftrightarrow E_d^{\mathrm{hol}}$).
  \item K.\ Costello and O.\ Gwilliam, \emph{Factorization Algebras
    in QFT} Vol~II, \S5.4 (chiral Lie structure on
    $\Omega^{0,\bullet}(\C^d, \g)$).
  \item Wave-13 G5: $\CE^*_{\dbar,\chir}(\cE_{\hCS}, \cO) = E_3$-algebra
    of observables; object-level input to Cycle~4.
  \item Vol~I Thm~A: chiral bar--cobar on curves; $d = 1$ boundary
    condition for Wave-14 J4's $d = 3$ lift.
\end{itemize}
